\setcounter{chapter}{1}
\setcounter{chapter}{1} 
\setcounter{table}{0}
\setcounter{figure}{0}
%\addcontentsline{toc}{chapter}{KHẢO SÁT VÀ PHÂN TÍCH HỆ THỐNG}
\chapter{KHẢO SÁT VÀ PHÂN TÍCH HỆ THỐNG}


\section{Khảo sát hiện trạng}
Sinh viên mô tả hiện trạng, quy trình nghiệp vụ hoặc bài toán thực tế liên quan đến đề tài.

\section{Phân tích yêu cầu}
Sinh viên trình bày yêu cầu chức năng, yêu cầu phi chức năng và các ràng buộc của hệ thống hoặc giải pháp.

\section{Phương pháp thực hiện}
Sinh viên mô tả phương pháp, quy trình, công cụ và dữ liệu sử dụng trong quá trình thực hiện đề tài.

\section{Kế hoạch thực hiện}
Sinh viên trình bày các giai đoạn công việc, mốc thời gian và sản phẩm dự kiến của từng giai đoạn.
Ví dụ, kế hoạch thực hiện có thể được trình bày như Bảng~\ref{tab:ke-hoach}.

\begin{table}[H]
\centering
\caption{Ví dụ bảng kế hoạch thực hiện}
\label{tab:ke-hoach}
\begin{tabular}{|c|p{5cm}|p{5cm}|}
\hline
\textbf{STT} & \textbf{Công việc} & \textbf{Kết quả dự kiến} \\
\hline
1 & Khảo sát và thu thập yêu cầu & Danh sách yêu cầu \\
\hline
2 & Phân tích và thiết kế & Tài liệu thiết kế \\
\hline
3 & Triển khai và kiểm thử & Sản phẩm và kết quả kiểm thử \\
\hline
\end{tabular}
\end{table}



\begin{figure}[htbp!]
 \begin{center}
\psfrag{z}{$\zeta$}
\psfrag{x}{$\xi$}
\psfrag{m}{$\xi_{\rm mid}$}
\psfrag{p}{$\xi_{\rm mid}^+$}
\psfrag{mm}{$\xi_{\rm mid}^-$}
\psfrag{xp}{$\xi_+$}
\psfrag{xm}{$\xi_-$}
\psfrag{xps}{$\xi'_+$}
\psfrag{xms}{$\xi'_-$}
\includegraphics[width=5.5cm]{\figpath refine.eps} 
\end{center}
\caption{Các tâm được thêm vào tại lân cận cạnh được đánh dấu. 
}
\label{refine}
\end{figure}
 Thuật toán trong \cite{DavyOanh11} làm mịn cạnh $\zeta \xi$ bằng cách chèn thêm tâm mới $\xi _ { \rm {mid } }$ là trung điểm của cạnh này. Còn trong \cite{ODP2017}, ngoài trung điểm còn thêm các điểm $ \xi _ { \rm{ mid } } ^ { + } , \xi _ { \rm{ mid } } ^ { - }$ nằm trên đường trung của cạnh $\zeta \xi$ và đối xứng nhau qua $\xi _ { \rm{mid }}$. Tuy nhiên các điểm này có thể nằm ở vị trí rất gần điểm đã tồn tại hoặc một điểm mới được tạo ra khi làm mịn một cạnh khác. Để tránh điều này các thuật toán trong  \cite{DavyOanh11, ODP2017} chỉ làm mịn nếu 
\begin{equation}\label{mu}
\dist(\xi',\Xi')\ge\mu\sep_{\xi'}(\Xi'),\quad\text{với $\mu\in(0,1)$,}
\end{equation}
tức là, việc chèn một tâm mới $\xi'$ trong tập hiện tại $\Xi'$ không làm giảm  
\emph{khoảng cách tách biệt địa phương} được xác định công thức
\begin{equation}\label{sep}
\sep_{\xi'}(\Xi'):=\frac{1}{4}\sum_{i=1}^4\dist(\xi_i,\Xi'\setminus\{\xi_i\}),
\end{equation}
trong đó $\xi_1,\ldots,\xi_4$ là bốn điểm gần nhất trong $\Xi'$  đến $\xi'$, và
$\dist(x,Y):=\inf\{\|x-y\|:\,y\in Y\}$
là khoảng cách từ điểm $x$ đến tập $Y$. 
% Tuy nhiên, chúng tôi đã kiểm tra 
% $\xi'\notin\partial\Omega$ không ở vị trí quá gần biên, 
% xem các Bước II.2 và II.3.i của Thuật toán~\ref{alg2Refi}.  

Nếu $\xi \in \partial \Xi$, thì tìm 2 điểm lân cận $\xi _ { - } , \xi _ { + } \in \partial \Xi$ và chèn thêm 2 trung điểm $\xi _ { - } ^ { \prime } , \xi _ { + } ^ { \prime }$ của các cạnh $\xi \xi _ { - }$ và $\xi \xi _ { + }$.



% 

\begin{algorithm}\label{alg1Refi}\rm {\bf Thuật toán làm mịn thích nghi  [Thuật toán 2, \cite{DavyOanh11}]} \\
\emph{Input:} Bộ tâm $\Xi$ và các khuôn điểm $\{\Xi_\zeta:\,\zeta\in\Xi_{int}\}$. \\
\emph{Output:} Bộ tâm được làm mịn $\Xi'$ và các giá stencil $\{ \Xi' _ { \zeta }: \zeta \in \Xi' \setminus \partial \Xi' \}$ . \\
\emph{Các tham số:}   $\gamma = 0.5$, $\mu = 0.7$. 
\begin{enumerate}[I.] 
\item\label{1step1} Tính độ lệch lớn nhất  $\bar\eps=\bar\eps(\Xi)$ và đánh dấu tất cả các cạnh $\zeta\xi, \xi \in \Xi_{\zeta}\setminus \{\zeta\}, \zeta \in \Xi_{\rm int}$ sao cho $\eps(\zeta, \xi) \ge \gamma\bar\eps$. Thiết lập  $\Xi':=\Xi$.
\item \label{1step2} \emph{Với mỗi} cạnh được đánh dấu $\zeta\xi$: 
\begin{enumerate}[1.]
\item Đặt $\xi':=(\zeta+\xi)/2$.
\item \emph{Nếu} $\dist\left( \xi' , \Xi' \right) \ge \mu \sep _ { \xi' } \left( \Xi' \right)$:
\begin{enumerate}[i.] 
\item Thiết lập  $\Xi': = \Xi' \cup \{ \xi'\}$.
\item \emph{Nếu} $\xi \in \partial \Xi$ thì thiết lập $\Xi': = \Xi' \cup \{ \xi' _ { - }, \xi' _ { + }\}$.
\end{enumerate}
\end{enumerate}
\item 
\begin{enumerate}[1.]
\item \emph{Với mỗi} $\zeta \in \Xi' \setminus \partial \Xi'$:
\begin{enumerate}[i.] 
\item Áp dụng Thuật toán \ref{alg1} để tìm $\Xi' _ { \zeta } \subset \Xi'$. 
\item Tính khoảng cách $d _ { \xi } = \| \zeta - \xi \| _ { 2 } , \xi \in \Xi' _ { \zeta } \setminus \{ \zeta \}$ và trung bình của chúng 
$d _ { a v } : = \frac { 1 } { \ell } \sum _ { \xi \in \Xi' _ { \zeta } \setminus \{ \zeta \} } d _ { \xi }$, trong đó $\ell = \# \Xi' _ { \zeta } - 1$ và 
đánh dấu tất cả các cạnh $\zeta \xi , \xi \in \Xi _ { \zeta } \setminus \{ \zeta \}$ sao cho $d _ { \xi } > 2 d _ { a v }$.
\end{enumerate}
\item Lặp lại Bước II đối với cạnh đã được đánh dấu trong Bước III.1 và tập tất cả các tâm mới đã được tạo trong tập $\Xi_{aux}$.
\item Áp dụng Thuật toán \ref{alg1} để tìm $\Xi' _ { \zeta }$ với mọi $\zeta \in \Xi_{aux}$ và tính lại $\Xi' _ { \zeta }$ với $\zeta \in \Xi' \setminus \Xi _ { a u x }$ sao cho $\dist  \left( \zeta , \Xi _ { a u x } \right) < d \left( \Xi' _ { \zeta }\right): = \max _ { \xi \in \Xi'_ { \zeta } } \| \zeta - \xi \| _ { 2 }$. 
\end{enumerate}
\end{enumerate}
\end{algorithm}
%  